\section{研究问题先于材料数量}

本章回答的投资问题是：在利润预告、行业景气和券商乐观判断同时出现时，哪些信息足以进入盈利与估值，哪些只能成为监控变量？核心判断是，恒逸石化已有足够的官方财务、项目、股本和外部原始研报支持周期归一化估值；但2026年分部利润、实际裂解价差、文莱税惠起止、员工计划费用和二期完整融资条款仍不完整。因此报告可以给出中性／观察与公允区间，却不能把高景气、二期成长或客户订单写成确定性价值。

证据治理并不是把所有来源平均对待。公司法定披露回答“发生了什么”，行业数据回答“外部环境如何”，券商报告提供“另一套预测与估值”，House模型负责把事实和假设转成可复算结论。层级越弱，允许承担的主张越窄；同一数字若来自不同口径，也不能因为数值相近就相互替代。

\begin{exhibitbox}[证据层级与允许用途]
\small
\noindent\begin{tabularx}{\textwidth}{L{2.25cm}L{2.6cm}L{2.8cm}L{2.6cm}X}
\toprule
\textbf{层级} & \textbf{代表来源} & \textbf{可证明} & \textbf{不可证明} & \textbf{估值用途} \\
\midrule
L1公司官方 & 年报、Q1、H1预告、转股、员工计划、二期公告 & 报告期财务、股本、产能、计划与公告边界 & 未披露分部利润、未来价差、最终认购与项目回报 & 主模型事实输入；预告标明未经审计 \\
L1公共机构 & IEA、EIA、统计局、郑商所、文莱政府 & 全球／行业供需、预测、方法与制度背景 & 恒逸实际采购、开工、客户或单吨利润 & 场景与方法，不直接进入公司EPS \\
L1原始券商PDF & 国信证券2026-07-04报告 & 券商预测、目标、方法与当时股本口径 & 公司指引或多券商一致预期 & 归一化后给20\%外部锚权重 \\
L2--L3行业代理 & 行业协会、商务预报转载 & PTA／聚酯开工、库存、加工费方向 & 恒逸装置和利润的一比一映射 & 验证情景，不作产能乘价差 \\
媒体摘要／搜索片段 & 7份券商转载摘要 & 市场叙事与可见字段 & 原始估值方法、完整预测和正式目标 & 权重为零 \\
House推断 & H2正常化、概率、PE与目标价 & 透明假设下的研究判断 & 公司承诺或确定收益 & 可复算，但必须给证伪条件 \\
\bottomrule
\end{tabularx}
\end{exhibitbox}
\sourcenote{来源登记、主张审计及官方／行业抓取清单，检索截止2026-07-22。}

这套层级直接影响估值：H1预告可以作为预测起点，但不能替代半年报；PTA加工费可以定义升级和降级阈值，但不能乘以2,150万吨参控股产能；国信目标价可以作为外部锚，但必须先把其约36.03亿股口径归一到38.80694189亿股，并同时处理员工计划现金。

\section{公司披露、派生数与预测必须分开}

同一页中最容易发生的错误，是把已披露、派生和预测混成一类。例如2026Q1毛利率12.8838\%是用收入减营业成本派生，数值可靠但不是公司单列口径；2026H1归母55--60亿元是公司预告，具有官方来源却未经审计；77.5亿元则完全是House预测。读者只有先看清类型，才能判断一个差异究竟来自经营、会计口径还是研究假设。

\begin{exhibitbox}[关键数字的状态标记]
\noindent\begin{tabularx}{\textwidth}{L{3.2cm}L{2.1cm}L{2.2cm}X}
\toprule
\textbf{指标} & \textbf{状态} & \textbf{数值} & \textbf{使用规则} \\
\midrule
2025A归母净利润 & 已审计披露 & 2.58亿元 & 历史低谷基线 \\
2026Q1归母／OCF & 官方季报 & 19.95／-25.73亿元 & 利润与现金质量的反差 \\
2026Q1毛利率 & 派生 & 12.8838\% & 以收入和营业成本复算，注明derived \\
2026H1归母预告 & 官方、未经审计 & 55--60亿元 & 预测起点；最终以半年报为准 \\
2026E House归母 & 研究预测 & 77.5亿元 & H1中值57.5 + H2正常化20.0 \\
2026E国信归母 & 外部券商预测 & 79.1亿元 & Street对照，不写成公司指引 \\
发布目标价 & House结论 & 15.7元 & 多锚模型；对应约4.3\%上行 \\
\bottomrule
\end{tabularx}
\end{exhibitbox}
\sourcenote{HY-OFF-2025AR；HY-OFF-2026Q1；HY-OFF-2026H1P；HY-BRK-GUOXIN；财务与模型截止2026-07-22。}

投资含义是，H1预告若最终落在区间内，只能证明预告可靠，不能自动证明H2和2027年也维持同一运行率。要升级估值，至少还需现金流、文莱分部、少数股东与国内库存共同支持；任何单一行业价差或券商摘要都不能替代这一组合证据。

\section{口径差异不是噪声，而是风险来源}

2025年报与2026Q1报告对期初资产和归母净资产存在重述差异；H1预告使用的去年同期同比基数也可能因同一控制下企业合并继续调整。这些差异不意味着披露错误，却会让同比增速、ROE和分部桥在未经对齐时失真。本报告因此优先使用绝对额与同口径期末值，并把预告同比只作方向说明。

债务口径也必须分开。2025年报总负债792.40亿元是资产负债表口径；评级报告刚性债务699.50亿元是机构重分类，其中短期刚性债务528.93亿元。前者用于负债率，后者用于偿债与利息压力分析。将两者直接相减或互称，会制造不存在的“其他债务”。同样，法定总股本、经济股本、转债if-converted股本和已披露稀释股本服务于不同问题，不能任选一个得到更好看的EPS。

\begin{sourcequalitybox}[正式边界]
本报告所有重要金额均尽量保留“期间、单位、状态和来源ID”。行业利用率不写成公司开工，区域运费不写成恒逸成本，券商预测不写成公司指引，需求锚不写成订单，规划产能不写成已投产销量。\textbf{投资上，任何跨越这些边界的新信息都必须先完成口径对齐，再决定是否调整盈利或目标价。}
\end{sourcequalitybox}

\section{主张审计：支持、推断与拒绝}

主张审计的目的，是让读者知道结论依赖什么，也知道报告主动拒绝了什么。H1利润、Q1现金流、少数股东损益、资产负债率和盘中价格属于已披露事实；“文莱是最大驱动”属于公司原因说明加支持性推断；“当前价反映多数基准利润”属于House模型反推；“20元以上为基准价值”“客户订单形成backlog”“北向资金看多”则没有证据支持。

\begin{exhibitbox}[高影响主张审计]
\small
\noindent\begin{tabularx}{\textwidth}{L{3.45cm}L{2.1cm}L{3.2cm}X}
\toprule
\textbf{主张} & \textbf{判定} & \textbf{报告采用表述} & \textbf{估值处理} \\
\midrule
H1归母55--60亿元 & 公司预告 & 未经审计，待半年报确认 & 情景起点 \\
H1主要来自核心经营 & 支持性推断 & 扣非接近归母，但缺分部拆分 & 排除大额非经常性主导，不精确拆项目 \\
文莱一期是主要驱动之一 & 公司原因+推断 & 满产满销和油品／芳烃盈利高 & 文莱周期敏感性，不反算单吨利润 \\
Q1现金转化弱 & 已披露 & 归母19.95亿元、OCF -25.73亿元 & 压低基准PE \\
当前价反映多数基准利润 & House推断 & 稀释敏感性隐含75.42亿元 & 中性评级依据 \\
20元以上是基准价值 & 拒绝 & 仅属牛情景尾部 & 不进入日常目标区间 \\
客户订单形成backlog & 未披露 & 商品销售不虚构订单能见度 & 权重为零 \\
北向／机构资金看多 & 未收集 & 不判断资金属性 & 权重为零 \\
\bottomrule
\end{tabularx}
\end{exhibitbox}
\sourcenote{data/claim\_audit.md，审计截止2026-07-22；HY-OFF-2026H1P；HY-OFF-2026Q1；HY-MKT-QUOTE。}

审计结果并不降低报告可用性，反而明确了更新路径：当半年报提供分部利润、现金流和子公司附注时，支持性推断可以升级为公司级证据；当员工计划过户与费用确认、转债进一步转股或二期贷款交割出现时，资本结构假设可以替换为实现值。没有这些新证据之前，结论保持有条件而非不断向乐观叙事漂移。

\section{Street证据为什么只有一份获得权重}

截至2026-07-22，案例归档八份去重卖方报告线索，其中只有国信证券2026-07-04报告取得原始PDF并完成哈希校验。其余七份只取得媒体转载摘要和受限下载页，虽然部分摘要显示盈利预测或目标价，但缺少完整估值方法、页码与正式文件，不能进入加权Street结果。

\begin{exhibitbox}[卖方证据覆盖与使用上限]
\noindent\begin{tabularx}{\textwidth}{L{3.1cm}R{1.7cm}R{1.8cm}L{2.2cm}X}
\toprule
\textbf{证据类别} & \textbf{数量} & \textbf{估值权重} & \textbf{代表} & \textbf{处理} \\
\midrule
原始PDF、字段完整 & 1 & 20\% & 国信证券 & 股本与员工计划现金归一后形成16.3130元Street锚 \\
媒体转载摘要 & 7 & 0 & 光大、中金、国海、中信建投 & 只观察盈利预测温度，不补齐缺失字段 \\
券商官方页／授权终端结构化导出 & 0 & 0 & 未取得 & 若后续取得，先归档复核再调整权重 \\
匿名聚合一致预期／搜索片段 & 不采用 & 0 & -- & 不承担目标价或盈利证明 \\
\bottomrule
\end{tabularx}
\end{exhibitbox}
\sourcenote{data/broker\_street\_consensus\_20260722.md；sources/broker-reports/2026-07-22/index.md；HY-BRK-GUOXIN，检索截止2026-07-22。}

只有一份正权重样本意味着Street锚不是“市场一致预期”，也不应得到比基本面更高的权重。未来若取得更多原始报告，新增样本可能提高也可能降低16.3130元锚；在此之前，报告宁可承认样本窄，也不把摘要数量伪装成共识深度。

\section{证据门禁对估值与风险的最终含义}

当前证据足以回答：恒逸具有什么资产、2025与2026Q1表现如何、H1预告区间是什么、股本和转债结构怎样、员工计划上限与价格如何、市场当时交易在什么位置。当前证据不足以回答：文莱2026年实际税后利润、各产品单吨盈利、二期融资成本、员工计划最终费用和其余库存股用途。

因此，估值上采用周期调整PE与三情景，而不是项目级伪精确SOTP或峰值DCF；风险上把现金流、杠杆和每股稀释与盈利风险并列；监控上把下一份半年报、三季报、员工计划过户、转债窗口和二期融资作为证据升级点。

\begin{keyinsight}[证据门禁的“所以呢”]
门禁状态为CONDITIONAL，不等于研究失效；它意味着目标价可以发布，但成长信用必须受限。\textbf{行动上，15.7元目标只依赖已披露资产、利润和资本结构，不为未交割二期融资、未披露订单、未知股份支付或峰值价差提前付费。}
\end{keyinsight}
